\documentclass[aspectratio=169,11pt]{beamer}
\usetheme[hideothersubsections=true]{Berkeley}

\usepackage[fontset=fandol]{ctex}
\usepackage{fontspec}
\usepackage{booktabs}
\usepackage{array}
\usepackage{tikz}
\usetikzlibrary{positioning}
\usepackage{listings}

\definecolor{UCBBlue}{HTML}{003262}
\definecolor{UCBGold}{HTML}{FDB515}
\definecolor{UCBFoundersRock}{HTML}{3B7EA1}
\definecolor{UCBMedalist}{HTML}{C4820E}
\definecolor{UCBGray}{HTML}{F5F5F5}

\setbeamertemplate{navigation symbols}{}

\setmainfont{Helvetica Neue}
\setsansfont{Helvetica Neue}
\setmonofont{Andale Mono}

\lstset{
  basicstyle=\ttfamily\small,
  keywordstyle=\color{UCBBlue},
  commentstyle=\color{UCBFoundersRock},
  breaklines=true,
  columns=fullflexible,
  frame=single,
  rulecolor=\color{UCBFoundersRock}
}

\title[LLM 辅助异步状态机跟踪]{一种基于 LLM 辅助的 Tokio 异步状态机动态跟踪分析}
\subtitle{从 GDB 原始日志到可解释报告}
\author{陈玺州}
\date{2026 年春夏开源操作系统训练营}

\begin{document}

\begin{frame}[plain]
  \titlepage
\end{frame}

\section{研究动机}
\subsection{问题意识}

\begin{frame}{问题意识}
  \begin{block}{目标}
    不只让异步爬虫能运行，而是解释 \texttt{async/await} 在运行时如何保存状态、返回 \texttt{Pending}、被 \texttt{Waker} 唤醒并继续执行。
  \end{block}

  \vspace{0.4em}
  \begin{itemize}
    \item 表层代码只是几个 \texttt{.await}
    \item 底层行为涉及 future frame、寄存器、栈帧、捕获变量区和 runtime 调度
    \item 核心挑战：把高层抽象还原为可观察的机器状态
  \end{itemize}
\end{frame}

\subsection{技术路线}

\begin{frame}{整体技术路径}
  \centering
  \begin{tikzpicture}[node distance=1.1cm, every node/.style={font=\large}]
    \node[draw=UCBBlue, thick, rounded corners, fill=UCBGray, minimum width=3.2cm, minimum height=1.0cm] (raw) {1. 获得原始数据};
    \node[draw=UCBFoundersRock, thick, rounded corners, fill=UCBGray, minimum width=3.2cm, minimum height=1.0cm, right=of raw] (ai) {2. LLM 辅助分析};
    \node[draw=UCBMedalist, thick, rounded corners, fill=UCBGray, minimum width=3.2cm, minimum height=1.0cm, right=of ai] (report) {3. 总结提炼};
    \draw[->, very thick, color=UCBBlue] (raw) -- (ai);
    \draw[->, very thick, color=UCBBlue] (ai) -- (report);
  \end{tikzpicture}

  \vspace{1.0em}
  \begin{alertblock}{核心优点}
    暴露 LLM 得出结果的路径：每个结论都可以回溯到命令、断点、日志、寄存器、内存地址和调用栈。
  \end{alertblock}
\end{frame}

\section{实验采集}
\subsection{原始数据}

\begin{frame}[fragile]{第一步：获得原始数据}
  \begin{lstlisting}[language=bash]
cargo build --bin async
timeout 35s rust-gdb -q target/debug/async -x gdb_scripts/gdb_track.gdb
  \end{lstlisting}

  \begin{itemize}
    \item 使用 dev profile，保留 debuginfo
    \item 先限制 \texttt{SPIDER\_CONCURRENCY=1}，观察单个状态机的完整生命周期
    \item 再追加 \texttt{SPIDER\_CONCURRENCY=4}，观察多个状态机实例并存
    \item 原始输出保存为 GDB 日志，再整理为中文分析版
  \end{itemize}
\end{frame}

\subsection{断点设计}

\begin{frame}{断点设计}
  \begin{tabular}{>{\ttfamily}p{0.43\linewidth}p{0.47\linewidth}}
    \toprule
    位置 & 观察目标 \\
    \midrule
    download\_async.rs:16 & 外层 create\_dir\_all.await \\
    download\_async.rs:25 & 每个 task 的 async block \\
    downloader\_async.rs:13--14 & reqwest send().await \\
    downloader\_async.rs:16--17 & response.bytes().await \\
    downloader\_async.rs:20 & 内部 create\_dir\_all(parent).await \\
    downloader\_async.rs:23 & tokio::fs::write(output, bytes).await \\
    \bottomrule
  \end{tabular}

  \vspace{0.8em}
  \begin{block}{设计原则}
    断点尽量卡在 await 前后，让日志呈现“保存状态 -> Pending -> 唤醒 -> 恢复”的切换过程。
  \end{block}
\end{frame}

\section{LLM 分析}
\subsection{分析流程}

\begin{frame}{第二步：AI 辅助分析}
  \begin{columns}[T,totalwidth=\textwidth]
    \begin{column}{0.32\textwidth}
      \begin{block}{结构化}
        按源码位置、状态机层级、寄存器和 backtrace 归类日志。
      \end{block}
    \end{column}
    \begin{column}{0.32\textwidth}
      \begin{block}{生成假设}
        基于 ABI、\texttt{Future::poll} 形态和内存内容推断地址含义。
      \end{block}
    \end{column}
    \begin{column}{0.32\textwidth}
      \begin{block}{证据对照}
        要求每个结论说明来自哪条日志、哪个地址或哪个变量。
      \end{block}
    \end{column}
  \end{columns}

  \vspace{0.8em}
  \begin{alertblock}{使用原则}
    LLM 不直接给最终裁决，而是帮助整理证据链、发现重复模式、标注不确定性。
  \end{alertblock}
\end{frame}

\subsection{证据链}

\begin{frame}{证据链示例}
  \begin{tabular}{>{\ttfamily}p{0.26\linewidth}p{0.62\linewidth}}
    \toprule
    观察值 & 解释依据 \\
    \midrule
    rdx = 0xc & 对应 "output/async" 的 12 字节长度 \\
    rcx = 0x636e...7475 & 小端序能读出 "ut/async" 片段 \\
    0x555556351ce0 & 附近内存包含 "output/async" \\
    0x7fffffffaff8 & backtrace 中多次作为外层 cx 出现 \\
    0x555556355d10 & per-task 捕获区，包含 task/url/output \\
    \bottomrule
  \end{tabular}

  \vspace{0.8em}
  \begin{block}{关键点}
    不是“AI 说它是什么”，而是“日志中的哪些事实支持这个判断”。
  \end{block}
\end{frame}

\subsection{人工校验}

\begin{frame}{第三步：人工提炼与校验}
  \begin{itemize}
    \item 稳定事实：反复出现、能和 backtrace 或内存内容相互印证
    \item 临时现象：源码行断点处的编译器临时值，不强行解释
    \item 结论边界：明确哪些地址能对应变量，哪些寄存器不能稳定解读
  \end{itemize}

  \vspace{0.6em}
  \begin{alertblock}{底层分析里的危险}
    LLM 容易把所有值都解释完整；但系统调试需要保留“不知道”，否则会把临时值包装成错误结论。
  \end{alertblock}
\end{frame}

\section{控制块视角}
\subsection{从状态机到控制块}

\begin{frame}{从状态机到控制块}
  \begin{columns}[T,totalwidth=\textwidth]
    \begin{column}{0.47\textwidth}
      \begin{block}{编译器生成的部分}
        \begin{itemize}
          \item \texttt{Future} frame
          \item 状态编号：\texttt{Unresumed}/\texttt{SuspendN}/完成
          \item 跨 \texttt{await} 仍需使用的局部变量
          \item 当前正在等待的子 future：\texttt{\_\_awaitee}
        \end{itemize}
      \end{block}
    \end{column}
    \begin{column}{0.47\textwidth}
      \begin{block}{运行时管理的部分}
        \begin{itemize}
          \item 是否 ready / running / complete
          \item \texttt{Waker} 和 ready queue 节点
          \item 任务归属的 executor / scheduler
          \item 完成结果和等待者
        \end{itemize}
      \end{block}
    \end{column}
  \end{columns}

  \vspace{0.7em}
  \begin{alertblock}{控制块视角}
    \texttt{Future} frame 解释“从哪里继续执行”；任务控制元数据解释“谁在什么时候再次 poll 它”。
  \end{alertblock}
\end{frame}

\subsection{控制块字段}

\begin{frame}{协程任务控制块包含什么}
  \small
  \begin{tabular}{>{\ttfamily}p{0.28\linewidth}p{0.60\linewidth}}
    \toprule
    字段类别 & 作用 \\
    \midrule
    state bits & 记录 running / notified / complete / cancelled 等生命周期状态 \\
    future frame ptr & 指向保存 async 状态机和跨 await 变量的对象 \\
    waker & I/O ready 后重新找到这个任务 \\
    queue links & 挂到 ready queue 或 \texttt{FuturesUnordered} 的内部队列 \\
    output / join & 任务完成后的返回值、错误和等待者关系 \\
    scheduler handle & 知道应该交给哪个 runtime / executor 继续调度 \\
    \bottomrule
  \end{tabular}

  \vspace{0.4em}
  \begin{block}{本实验里的对应关系}
    每个下载并不是显式 \texttt{tokio::spawn} 出来的 \texttt{JoinHandle}；更直接的调度单位是 \texttt{buffer\_unordered} / \texttt{FuturesUnordered} 内部保存的子 future 节点。
  \end{block}
\end{frame}

\subsection{Waker 与唤醒}

\begin{frame}[fragile]{Waker：控制块的回指针}
  \begin{lstlisting}
RawWaker {
  data:   *const (),          // 指向 task node / 控制块
  vtable: *const RawWakerVTable
}

RawWakerVTable:
  clone       -> clone_arc_raw
  wake        -> wake_arc_raw
  wake_by_ref -> wake_by_ref_arc_raw
  drop        -> drop_arc_raw
  \end{lstlisting}

  \vspace{0.4em}
  \begin{itemize}
    \item 子 future 返回 \texttt{Pending} 时，通过 \texttt{Context} 注册 \texttt{Waker}
    \item I/O ready 后调用 \texttt{wake\_by\_ref}，把对应节点重新标记为 ready
    \item 下一轮 \texttt{FuturesUnordered::poll\_next} 选择 ready 子 future 继续 poll
  \end{itemize}
\end{frame}

\subsection{与 OS TCB 对比}

\begin{frame}{与操作系统 TCB 的对照}
  \begin{tabular}{p{0.28\linewidth}p{0.30\linewidth}p{0.32\linewidth}}
    \toprule
    维度 & OS thread TCB & Tokio / Future 任务节点 \\
    \midrule
    调度单位 & 内核线程 & 用户态 future 实例 \\
    保存现场 & 寄存器 + 独立线程栈 & future frame + 状态位 \\
    等待 I/O & 阻塞线程，内核切换 & 返回 \texttt{Pending}，注册 \texttt{Waker} \\
    恢复方式 & 调度回同一线程上下文 & executor 再次 \texttt{poll} \\
    局部变量 & 主要在线程栈中 & 跨 \texttt{await} 的变量进入 frame \\
    \bottomrule
  \end{tabular}

  \vspace{0.8em}
  \begin{alertblock}{关键差异}
    Tokio task 不保存完整 CPU 上下文和私有 OS 栈；一次 \texttt{poll} 临时借用 worker 栈，持久状态保存在 future frame 和任务元数据中。
  \end{alertblock}
\end{frame}

\section{状态机结论}
\subsection{嵌套结构}

\begin{frame}[fragile]{核心结论：嵌套状态机}
  \begin{lstlisting}
外层下载调度状态机
  -> collect / buffer_unordered / FuturesUnordered
  -> per-task async move 状态机
  -> CrawlTask::run_async_to 状态机
  -> downloader_async 状态机
  -> reqwest send / bytes / tokio fs 子状态机
  \end{lstlisting}

  \begin{itemize}
    \item 每个 await 把后续要用的局部变量保存到 future frame
    \item 子 future Pending 后，控制权返回 runtime
    \item I/O ready 后，Waker 触发 runtime 再次 poll
    \item 状态机跳回 await 后的源码位置继续执行
  \end{itemize}
\end{frame}

\subsection{单并发}

\begin{frame}{单并发观察}
  \begin{itemize}
    \item line 16 $\rightarrow$ line 17：外层 create\_dir\_all.await 返回
    \item line 25 反复命中：同一个 per-task future 被多次 poll
    \item line 13--14 $\rightarrow$ line 15：send().await 返回 Response
    \item line 16--17 $\rightarrow$ line 19：bytes().await 返回 Bytes
    \item line 23：write(output, bytes).await 使用已保存的状态
  \end{itemize}

  \begin{block}{结论}
    单并发适合观察一个状态机的保存、Pending、唤醒和恢复。
  \end{block}
\end{frame}

\subsection{并发 4}

\begin{frame}{并发 4 观察}
  \begin{tabular}{>{\raggedright\arraybackslash}p{0.30\linewidth}>{\ttfamily}p{0.42\linewidth}}
    \toprule
    任务 & 捕获变量区地址 \\
    \midrule
    北京大学 & 0x555556355d10 \\
    清华大学 & 0x555556353f10 \\
    中国人民大学 & 0x555556353b00 \\
    北京师范大学 & 0x555556340590 \\
    \bottomrule
  \end{tabular}

  \vspace{0.8em}
  \begin{itemize}
    \item 同一种 async block 类型可以同时有多个实例
    \item 后续日志中 \texttt{rsi} 在四个地址之间轮换
    \item \texttt{FuturesUnordered} 保存子 future，并借助 waker/ready queue 调度
  \end{itemize}
\end{frame}

\section{方法反思}
\subsection{可解释性}

\begin{frame}{可解释性来自哪里}
  \begin{enumerate}
    \item 命令可复现：构建命令、GDB 脚本、并发参数明确
    \item 证据可追踪：结论能回到断点、寄存器、内存和 backtrace
    \item LLM 路径可暴露：结构化、假设、证据对照、人工校验分层清楚
    \item 不确定性可保留：不能稳定对应的值不写成确定结论
  \end{enumerate}
\end{frame}

\subsection{LLM 使用}

\begin{frame}{如何更好地使用 LLM}
  \begin{itemize}
    \item 先给原始数据，再要求分析
    \item 把大问题拆成可验证的小问题
    \item 要求输出证据链，而不是只输出结论
    \item 主动标注不确定性和反例
    \item 最终报告由人收束，保证问题意识和结论边界
  \end{itemize}

  \vspace{0.7em}
  \begin{alertblock}{一句话}
    LLM 的价值不是替代推理，而是把推理过程显性化、结构化、可复查。
  \end{alertblock}
\end{frame}

\subsection{总结}

\begin{frame}{总结}
  \begin{block}{技术收获}
    看清了 Tokio 异步程序中“一堆状态机”在内存和调度上的实际表现。
  \end{block}

  \begin{block}{方法收获}
    建立了从原始日志到 LLM 辅助分析，再到人工提炼报告的可解释工作流。
  \end{block}

  \begin{alertblock}{训练营后的理解}
    系统软件里真正有价值的结论，不是听起来合理，而是能被日志、代码和机器状态共同支撑。
  \end{alertblock}
\end{frame}

\end{document}
